\begin{abstract}
As LLM-based agents operate over sequential multi-step reasoning, hallucinations arising at intermediate steps risk propagating along the trajectory, thus degrading overall reliability.
Unlike hallucination detection in single-turn responses, diagnosing hallucinations in multi-step workflows requires identifying which step causes the initial divergence. To fill this gap, we propose a new research task, \textbf{automated hallucination attribution} of LLM-based agents, aiming to identify the step responsible for the hallucination and explain why. To support this task, we introduce AgentHallu, a comprehensive benchmark with: (1) 693 high-quality trajectories spanning 7 agent frameworks and 5 domains, (2) a hallucination taxonomy organized into 5 categories (Planning, Retrieval, Reasoning, Human-Interaction, and Tool-Use) and 14 sub-categories, and (3) multi-level annotations curated by humans, covering binary labels, hallucination-responsible steps, and causal explanations.
We evaluate 13 leading models, and results show the task is challenging even for top-tier models (like GPT-5, Gemini-2.5-Pro). 
The best-performing model achieves only 41.1\% step localization accuracy, where tool-use hallucinations are the most challenging at just 11.6\%.
We believe AgentHallu will catalyze future research into developing robust, transparent, and reliable agentic systems.
\end{abstract}